% ==========================================================
% 000_frontmatter.tex
% FRONTMATTER (FAS Technical Note)
% ==========================================================

\begin{abstract}
Forecasting pipelines often assume that all slices of demand are equally amenable to learning,
evaluation, and control once sufficient data and model sophistication are applied. In operational
settings, however, this assumption frequently fails. Certain demand slices exhibit structural
pathologies—such as extreme intermittency, unstable tail behavior, or recurrent spike-driven
error—that render both model fitting and downstream readiness control unreliable, regardless of
model class. Attempting to forecast such slices can introduce spurious confidence, inflate
computational cost, and obscure genuine signal elsewhere in the system.

This technical note introduces the \emph{Forecast Admissibility Surface (FAS)}, a deterministic,
slice-keyed construct that partitions a forecast domain into admissible regions for modeling and
control. FAS operates on baseline-derived error anatomy—such as spike frequency, tail magnitude,
and support size—to classify each slice as \Allowed, \Conditional, or \Blocked prior to advanced
model training. FAS is not a performance metric and does not select or tune forecasting models.
Instead, it serves as an auditable pre-model gating mechanism that constrains where learning,
evaluation, and readiness policies may be meaningfully applied within the Forecast Readiness
Framework.
\end{abstract}
